\documentclass[UTF8,fontset=fandol,11pt,a4paper]{ctexart}
\usepackage[a4paper,margin=1.75cm,headheight=15pt]{geometry}
\usepackage{amsmath,amssymb,mathtools,booktabs,tabularx,array,enumitem}
\usepackage[most]{tcolorbox}
\usepackage{xcolor,fancyhdr,hyperref,microtype,lastpage}
\newcolumntype{Y}{>{\raggedright\arraybackslash}X}
\newcolumntype{L}[1]{>{\raggedright\arraybackslash}p{#1}}
\newcolumntype{C}[1]{>{\centering\arraybackslash}p{#1}}
\setlength{\parindent}{2em}
\setlength{\parskip}{0.34em}
\setlength{\emergencystretch}{3em}
\renewcommand{\arraystretch}{1.2}
\setlength{\tabcolsep}{3pt}
\setlist[itemize]{itemsep=0.18em,topsep=0.35em,leftmargin=2em}
\setlist[enumerate]{itemsep=0.18em,topsep=0.35em,leftmargin=2em}
\definecolor{deep}{RGB}{38,58,72}
\definecolor{soft}{RGB}{244,247,248}
\definecolor{greenish}{RGB}{52,91,74}
\definecolor{warnc}{RGB}{136,61,58}
\definecolor{purple}{RGB}{84,72,112}
\hypersetup{colorlinks=true,linkcolor=deep,urlcolor=purple,pdftitle={CO-01 数据表示与运算}}
\pagestyle{fancy}
\fancyhf{}
\lhead{\small\color{deep}\textbf{408 · 计算机组成原理 CO-01}}
\rhead{\small\color{deep}Bits · Meaning · Operation · Loss}
\cfoot{\small 第 \thepage\ 页 / 共 \pageref{LastPage} 页}
\renewcommand{\headrulewidth}{0.4pt}
\newtcolorbox{corebox}[1]{breakable,colback=soft,colframe=deep,title=\textbf{#1},boxrule=0.8pt,arc=2mm}
\newtcolorbox{methodbox}[1]{breakable,colback=white,colframe=greenish,title=\textbf{#1},boxrule=0.7pt,arc=1.5mm}
\newtcolorbox{warnbox}[1]{breakable,colback=white,colframe=warnc,title=\textbf{#1},boxrule=0.7pt,arc=1.5mm}
\newcommand{\kw}[1]{\textbf{\color{deep}#1}}
\title{\vspace{-1.1cm}\textbf{数据表示与运算}\\[0.4em]\large 有限位宽怎样保持可解释，并显式暴露信息丢失}
\author{408 · 计算机组成原理 Topic CO-01}
\date{}

\begin{document}
\maketitle

\begin{corebox}{母问题：同一串比特怎样得到确定的数学与硬件结果？}
比特没有自带 signed、float 或 address 语义。完整推理链是
\[
\begin{aligned}
&\text{Bit Pattern}\to\text{Interpretation}\to\text{Exact Value}\\
&\to\text{Operation}\to\text{Finite-width Result}\to\text{Lost Information / Flags}
\end{aligned}
\]
任何答案都应先固定 $(\text{width},\text{interpretation},\text{operation},\text{destination})$；溢出、舍入和异常都发生在“精确数学对象重新进入有限可表示集合”的边界。
\end{corebox}

\begin{methodbox}{本册定位与 Owner 边界}
\begin{itemize}
\item \kw{Owns：}位权、定长整数解释、原码/补码、扩展/截断/移位、carry/overflow、乘除基本算法、目标位宽判溢出、IEEE 754 分类/运算/舍入。
\item \kw{Uses：}题设或 ISA 给出的结果宽度、异常语义、舍入模式和状态位定义。
\item \kw{Stops：}具体指令编码和软件可见异常归 CO-02，完整 ALU/控制通路归 CO-03，高级语言规则归机器级映射接口。
\end{itemize}
\end{methodbox}

\tableofcontents
\newpage

\section{表示系统：位串、解释与可表示集合}
一个 $n$ 位位串只是集合 $\{0,1\}^n$ 的元素。unsigned 把它解释为
\[
U(b)=\sum_{i=0}^{n-1}b_i2^i,\qquad 0\le U\le 2^n-1;
\]
补码把最高位赋负权：
\[
T(b)=-b_{n-1}2^{n-1}+\sum_{i=0}^{n-2}b_i2^i,
\quad -2^{n-1}\le T\le2^{n-1}-1.
\]
同一位串可同时作为无符号数、有符号数、掩码、地址或浮点字段；意义来自消费它的操作，而不是某一位的外观。

原码把符号与绝对值分开，存在 $+0/-0$，加减需额外处理符号和大小。补码把定长加减统一到模 $2^n$ 环，只有一个零，代价是范围不对称：最小负数没有同宽正数。

\section{定长加减：同一低位结果，不同溢出解释}
$n$ 位加法器保留
\[
r=(a+b)\bmod 2^n.
\]
unsigned overflow 关注最高位 carry；signed overflow 关注精确结果是否落在补码范围。加法中“同号输入得到异号输出”可判 signed overflow；等价电路判据是符号位的进位与出位异或。

\begin{center}\small
\begin{tabularx}{\linewidth}{L{3.0cm}YY}\toprule
4 位运算 & unsigned 解释 & 补码解释\\\midrule
$1111+0001\to0000$ & $15+1$ 越界，carry=1 & $-1+1=0$，无 overflow\\
$0111+0001\to1000$ & $7+1=8$，无 unsigned 越界 & $7+1$ 越界，overflow=1\\
$1000$ 取负仍为 $1000$ & 位模式变换结果 & $-8$ 的同宽正数不可表示\\
\bottomrule
\end{tabularx}
\end{center}

减法先明确硬件采用 $a+(\overline b+1)$，再按题目对 carry/borrow 的定义解释。零、负、carry、overflow 是不同谓词；不能由一个标志推出另一个解释。

\section{扩展、截断与移位都在选择保留什么}
zero extension 保留 unsigned 值；sign extension 保留补码值。截断到低 $k$ 位相当于模 $2^k$，通常不保留数学值。判断窄化是否无损的通用办法是：截断后按原解释扩回，是否恢复原位串/数值。

逻辑移位补零；算术右移复制补码符号位，常对应向负无穷取整的除 $2^k$，不自动等于高级语言的有符号除法。左移低位结果与乘 $2^k$ 同余，但是否溢出仍取决于目标解释和被丢高位。

\begin{methodbox}{不变量}
每次扩展、截断或移位都问：要保留的是\kw{位模式、无符号值、补码值，还是仅保留模结果}？若答案不明确，后续标志判断没有语义基础。
\end{methodbox}

\section{乘法：完整乘积、部分积与目标宽度}
两个 $n$ 位无符号数的完整乘积最多需要 $2n$ 位；两个 $n$ 位补码数的完整乘积也用 $2n$ 位承载。若只保留低 $n$ 位，unsigned 与补码乘法的低位位串相同，但高位和溢出解释不同。

窄结果判据：unsigned 要求完整乘积高 $n$ 位全零；补码要求完整乘积高 $n$ 位等于低半部符号位的扩展。更一般地，先得到足宽精确积，再检查是否能无损压回目标格式。

\subsection{移位加法与原码一位乘}
无符号乘法把乘数每一位 $q_i$ 选择的 $M\cdot2^i$ 累加。序贯乘法器用部分积、乘数和被乘数移位，以时间换加法器/寄存器资源；并行阵列同时生成和压缩部分积，以面积与关键路径换吞吐。

原码一位乘先取符号异或，再对绝对值做无符号部分积；它不是补码 Booth 算法的同义词。题目给出联合寄存器时，必须先写每段宽度、移位方向和每拍不变量，不能跨教材套门数或节拍数。

\subsection{Booth 的生成逻辑}
Radix-2 Booth 用相邻位 $(Q_0,Q_{-1})$ 压缩连续的 1：$01$ 加 $M$，$10$ 减 $M$，$00/11$ 不变，然后对联合寄存器算术右移。它利用
\[
00111100_2=01000000_2-00000100_2
\]
把一串加法变为边界处加/减。若题目使用更高基 Booth、不同位对次序或不同寄存器布局，先声明版本；“Booth 对所有输入都更快”不是稳定结论。

\section{除法：商余关系是最终不变量}
除法的稳定正确性条件是
\[
X=YQ+R,\qquad |R|<|Y|,
\]
余数符号和商的取整方向服从题设/ISA。恢复余数法在试减为负时恢复；不恢复余数法让当前余数符号决定下一拍加还是减，省去立即恢复，但结束时可能需要校正。

算法题先固定 partial remainder、quotient/dividend、divisor 和附加位布局。每拍写判断位、加减对象、移位方向和上商位；最终用商余关系回代。原码与补码算法、左移余数与右移除数版本不能混用同一口诀。

除数为零可由零检测逻辑发现，但架构结果可能是同步异常、规定值或其他语义。补码最小负数除以 $-1$ 会产生超出同宽 signed 范围的商；“同位宽整数除法不溢出”是错误命题。

\section{浮点：在非均匀格点上表示实数近似}
普通二进制规格化数为
\[
x=(-1)^s(1.f)_2\,2^{E-Bias}.
\]
符号选方向，阶码选尺度，尾数选该尺度上的格点。阶码全 0 或全 1 必须分支处理 zero、subnormal、infinity 和 NaN，不能套普通隐藏位公式。

subnormal 使用 $0.f$ 和最小正常指数，提供渐进下溢；infinity 与 NaN 的传播、比较和异常细节以标准/题设为准。把整个 IEEE 754 编码称作“原码”会掩盖偏置阶码、隐藏位和特殊编码。

相邻可表示数的绝对间距随指数增长，因此“大数精度更高”与“有效位数固定”不是一回事：有效相对精度近似稳定，但绝对间距变大。多数十进制小数在二进制中无限展开，输入编码时已经发生第一次舍入。

\section{浮点运算链与舍入证据}
加减法的生成链是
\[
\text{Classify}\to\text{Align}\to\text{Significand Operation}\to\text{Normalize}\to\text{Round}\to\text{Encode/Signal}.
\]
对阶右移较小尾数，sticky 位应记录所有被移出低位的 OR；guard、round、sticky 与保留末位共同决定舍入。round-to-nearest, ties-to-even 下，恰好中点时让保留结果最低位为偶数；其他舍入模式必须另算。

乘除先处理符号和阶码，再算有效数，规格化、舍入并检查 exponent range。overflow 与 underflow 是目标格式边界；inexact 表示精确结果不在目标格点上。只看阶差达到尾数位数就断言小数完全无影响，可能漏掉 GRS 对最终舍入的作用。

\begin{warnbox}{三个浮点反例}
$0.1+0.2$ 涉及输入编码、运算和结果编码的多次有限映射；小量对阶后主体尾数为零，sticky 仍可能改变舍入；浮点加法不满足结合律，括号会改变中间舍入点。
\end{warnbox}

\section{五列概念边界}
\begin{center}\small
\begin{tabularx}{\linewidth}{L{1.8cm}C{0.35cm}L{1.9cm}YY}\toprule
概念 A & $\neq$ & 概念 B & 真正区别与题面信号 & 混淆后果\\\midrule
bit pattern & $\neq$ & value & 存储对象 vs 解释后的数学对象 & 同位串误判\\
carry & $\neq$ & overflow & unsigned 出位 vs signed 越界 & 标志互推\\
sign extension & $\neq$ & zero extension & 保补码值 vs 保 unsigned 值 & 扩宽变值\\
full product & $\neq$ & narrow result & 足宽精确积 vs 目标寄存器保留部分 & 漏乘法溢出\\
divide-by-zero detect & $\neq$ & ISA exception & 电路条件 vs 软件可见契约 & 从门电路推出 OS 行为\\
rounding & $\neq$ & truncation & 映射到最近/方向格点 vs 直接丢低位 & 浮点结果错一 ULP\\
\bottomrule
\end{tabularx}
\end{center}

\section{做题控制协议}
\begin{enumerate}
\item 写位宽、解释、运算和目标格式；
\item 先写可表示集合与精确数学结果；
\item 用足够中间位执行扩展、加减、部分积或商余更新；
\item 明确第一次丢信息的位置与舍入/截断规则；
\item 分开判断 carry、overflow、zero、sign、inexact 等状态；
\item 乘除用积回算或 $X=YQ+R$；浮点用分类、数量级和反向解码校验；
\item 最后才套 ISA 对标志、异常、饱和或返回值的契约。
\end{enumerate}

\begin{methodbox}{压缩成第一动作}
看到表示或运算题，先问：\kw{这串位按什么格式解释？精确结果在哪一步被压回多少位，丢掉了什么证据？}
\end{methodbox}

\section{全科坐标与跨册交接}
\begin{corebox}{Map Coordinate：ISA State / Semantics}
CO-01 负责把无语义的位串解释成给定位宽下的值、结果和信息损失证据。它向 CO-02/03 输出可消费的 typed value、宽度、符号扩展/舍入结果与标志候选；异常是否成为软件可见事件仍由 ISA 契约决定。
\end{corebox}
本册的局部成本包括运算位宽、加法器/乘除迭代次数、关键路径、舍入逻辑和精度损失。它不直接给出程序总时间，性能题要把这些参数交给 `IC x CPI x clock period` 的全局成本口径。

\section{考纲映射与来源处理}
\begin{center}\small
\begin{tabularx}{\linewidth}{L{3.0cm}YY}\toprule
考点 & 本册机制 & 接口\\\midrule
整数表示与运算 & 位权、补码、扩展、标志与目标位宽 & CO-02/03\\
乘法器与 Booth & 部分积、联合寄存器、相邻位重编码 & CO-03 Use\\
除法与溢出 & 商余不变量、恢复/不恢复、边界检测 & ISA exception\\
IEEE 754 & 分类、对阶、规格化、GRS 舍入、特殊值 & ISA rounding/status\\
\bottomrule
\end{tabularx}
\end{center}

\begin{corebox}{一句话}
有限硬件的核心不是“把数学算错”，而是按明确格式把无限或足宽结果投影回有限集合；正确答案必须同时给出保留结果、丢失信息和其软件可见解释。
\end{corebox}

本册吸收旧《定点小数》《IEEE》《浮点数》及十二张定点/浮点/乘除卡片；并行阵列只保留空间换时间的 Extension，具体 Booth/不恢复余数变体服从题设，空白《运算电路》不产生知识更新。

\end{document}
